\scanpage{17}
% Foliotation archivistique au crayon, en bas à gauche : 16.
le sous-corps (ou mieux, sous-extension) de $\overline{\mathbb Q}$ qu'ils engendrent, ext. finie de $\mathbb Q$. On voit de suite qu'il est indépendant du choix de $s$, et que c'est la plus petite sous-extension de $\overline{\mathbb Q}/\mathbb Q$ telle que $\widehat{\mathcal X}$ puisse se descendre \inserted{en $\widehat X_K$} à $K$ de façon que les él. de $S$ deviennent rat$/K$.
De plus la descente en question est unique, à isom unique près. Les autom \inserted{géom} de $(\widehat{\mathcal X},S)$ (ou encore de $\mathcal X$) forment un groupe fini ($\hookrightarrow\unsure{\mathfrak S_S|\mathbb G_S}$), et ce groupe opère \unsure{tt. déjà|tt entier} sur $\widehat X_K$.
